\section{Others}
\subsection{Game Theory \texorpdfstring{\&}{&} Nim}

\textbf{Standard Nim}: XOR sum of piles $= 0 \implies$ Second player wins. $\neq 0 \implies$ First player wins.

\textbf{Subtraction Game}: Can only take up to $k$ stones. XOR sum of $(pile_i \bmod (k+1))$.

\textbf{Index-k Nim}: Can remove from up to $k$ piles at once. Write sizes in binary. Sum each column. If all column sums $\equiv 0 \pmod{k+1}$, second player wins.

\textbf{Mis\`{e}re Nim} (Last player to move LOSES): If all piles have size 1, Second wins if $N$ is ODD. Otherwise, play like Standard Nim (XOR sum).

\subsection{SOS DP}
\begin{minted}{c++}
void zeta_subset(vector<int64_t> &dp, int bits) {
  for (int b = 0; b < bits; b++)
    for (int mask = 0; mask < (1 << bits); mask++)
      if (mask >> b & 1)
        dp[mask] += dp[mask ^ (1 << b)];
}
void zeta_superset(vector<int64_t> &dp, int bits) {
  for (int b = 0; b < bits; b++)
    for (int mask = 0; mask < (1 << bits); mask++)
      if (!(mask >> b & 1))
        dp[mask] += dp[mask | (1 << b)];
}
\end{minted}
\begin{minted}{c++}
void Solve() {
  ll n; cin >> n;
  vector<ll> arr(n);
  for (auto& x : arr) cin >> x;
  ll N = 1 + __lg(*max_element(arr.begin(),
  arr.end()));
  vector<ll> sos(1 << N, 0), sos2(1 << N, 0);
  for (auto it : arr) {
    sos[it]++;
    sos2[((1 << N) - 1) ^ it]++;
  }
  for (int i = 0; i < N; ++i)
    for (int mask = 0; mask < (1 << N); ++mask)
      if (mask & (1 << i))
        sos[mask] += sos[mask ^ (1 << i)],
        sos2[mask] += sos2[mask ^ (1 << i)];
  for (auto& it : arr) {
    // count y: x|y == x
    cout << sos[it] << ' ';
    // count y: x&y == x
    cout << sos2[((1 << N) - 1) ^ it] << ' ';
    // count y: x&y != 0
    cout << n - sos[((1 << N) - 1) ^ it] << ' ';
  }
}
\end{minted}
\subsection{Sum of Digits up to N}
\begin{minted}{c++}
uint64_t sumDigitsUpTo(uint64_t n) {
  if (n <= 0) return 0;
  uint64_t sum = 0;
  for (uint64_t f = 1; f <= n; f *= 10) {
    uint64_t hi = n / (f * 10);
    uint64_t cur = (n / f) % 10;
    uint64_t lo = n % f;
    sum += hi * 45LL * f;
    sum += cur * (lo + 1);
    sum += f * (cur * (cur - 1) / 2);
  }
  return sum;
}
\end{minted}
\subsection{Number of Digits up to N}
\begin{minted}{c++}
int64_t get(int64_t n) {
  int64_t base = 1, digits = 1, now = 0;
  while (base * 10 <= n) {
    base *= 10;
    now += (base - (base / 10 == 0 ? 1 :
           base / 10)) * digits;
    digits += 1;
  }
  n -= (base - 1);
  now += n * digits;
  return now;
}
\end{minted}
\subsection{Subset Sum (Bitset)}
Works when sum of values $\le$ N. TC: O(N/32 * sqrt(N))
\begin{minted}{c++}
const int N = 1e6 + 6;
bitset<N> dp;
dp.set(0);
for (auto &[w, c] : comp) {
  for (int i = 0; 1 << i <= c; ++i) {
    dp |= dp << (w * (1 << i));
    c -= 1 << i;
  }
  dp |= dp << (w * c);
}
dp.test(k) // if sum k is possible
\end{minted}

OR convolution can be done using SOS DP.

\subsection{All Submasks in $O(3^N)$}
\begin{minted}{c++}
for (int mask = 0; mask < 1 << n; mask++) {
  for (int sub = mask; sub > 0;
       sub = (sub - 1) & mask) {
  }
}
\end{minted}
